\section{Vorteile, Nachteile und geeignete Einsatzbedingungen}\label{sec:bewertung}
\textbf{Vorteile.} Der Verifier benötigt kein gemeinsames langfristiges Geheimnis. Zufällige private Schlüssel vermeiden schwache Benutzerpasswörter; öffentliche Schlüssel erlauben unter der DLP-Annahme keine effiziente Schlüsselrückgewinnung. Rechen- und Kommunikationsaufwand sind gering, Commitments lassen sich vorberechnen. Simulierbare Transkripte liefern allein keinen unabhängigen Signaturnachweis gegenüber Dritten.

\textbf{Grenzen.} Erforderlich sind Interaktion, vertrauenswürdige Schlüsselregistrierung, sichere Zufallswerte und korrekte Validierung. Das Verfahren bietet weder Verschlüsselung noch automatisch gegenseitige Authentifizierung oder Schutz vor Live-Weiterleitung. Es ist nicht postquantensicher: Shors Algorithmus löst das DLP auf einem ausreichend leistungsfähigen fehlertoleranten Quantencomputer effizient \cite{shor1997}.

\textbf{Eignung.} Schnorr lohnt sich in spezialisierten Protokollen, die Schlüsselkenntnis prüfen und bereits sichere Schlüsselverwaltung sowie eine Bindung an den Sitzungskontext gewährleisten. Öffentlich überprüfbare Nachrichtenbestätigungen erfordern Signaturen; nichtinteraktive oder postquantenresistente Nachweise benötigen erweiterte beziehungsweise andere Konstruktionen.
